\begin{thebibliography}{99}

\bibitem{interior_votos} 
Ministerio del Interior (2024). \textit{Área de descargas de procesos electorales}. Gobierno de España. Recuperado de \url{https://infoelectoral.interior.gob.es/es/elecciones-celebradas/area-de-descargas/}

\bibitem{ine_renta} 
Instituto Nacional de Estadística (2024). \textit{Atlas de distribución de renta de los hogares}. Serie 2015-2023. Recuperado de \url{https://www.ine.es/dynt3/inebase/index.htm?padre=12385&capsel=12384}

\bibitem{nastat_navarra} 
Instituto de Estadística de Navarra (2024). \textit{Estadística de renta de las personas físicas y familia}. Recuperado de \url{https://nastat.navarra.es/es/tablas_powerbi/-/tag/estadistica-renta}

\bibitem{sepe_empleo} 
Servicio Público de Empleo Estatal (2024). \textit{Datos estadísticos de paro y contratos por municipios}. Recuperado de \url{https://www.sepe.es/HomeSepe/que-es-el-sepe/estadisticas/datos-estadisticos/municipios.html}

\bibitem{ine_padron} 
Instituto Nacional de Estadística (2024). \textit{Padrón municipal de habitantes}. Recuperado de \url{https://www.ine.es/dyngs/INEbase/categoria.htm?c=Estadistica_P&cid=1254734710990}

\bibitem{ign_geo} 
Instituto Geográfico Nacional (2024). \textit{Delimitaciones Territoriales y Cartografía Municipal (NGMEP)}. Centro de Descargas del CNIG. Recuperado de \url{https://centrodedescargas.cnig.es/CentroDescargas/detalleArchivo?sec=9000004#}

\bibitem{cis_ideologia} 
Centro de Investigaciones Sociológicas (2022). \textit{Barómetro de diciembre de 2022 (Estudio 3388)}. Recuperado de \url{https://www.cis.es/}

% --- Nuevas referencias del Marco Teórico ---

\bibitem{alaminos2023} 
Alaminos-Fernández, A. F., \& Alaminos, A. (2023). \textit{Métodos y Modelos para la Predicción Electoral. Una Guía Práctica}. Universidad de Alicante.

\bibitem{albaladejo2025} 
Albaladejo Gutiérrez, G. (2025). \textit{Análisis predictivo del comportamiento electoral en España}. Trabajo de Fin de Grado, Universidad Pontificia Comillas.

\bibitem{deb2024} 
Deb, S., Roy, A., \& Das, S. (2024). \textit{Forecasting elections from partial information using a Bayesian model}. Journal of Applied Statistics, 51(2), 310-328.

\bibitem{kennedy2017} 
Kennedy, R., Wojcik, S., \& Lazer, D. (2017). \textit{Improving election prediction internationally}. Science, 355(6324), 515-520.

\bibitem{lewisbeck2019} 
Lewis-Beck, M. S., \& Stegmaier, M. (2019). \textit{Economic Voting}. The SAGE Handbook of Electoral Behaviour.

\bibitem{lipset1959} 
Lipset, S. M. (1959). \textit{Some social requisites of democracy: Economic development and political legitimacy}. American Political Science Review, 53(1), 69-105.

\bibitem{pavia2021} 
Pavía, J. M., \& Cantarino, I. (2021). \textit{Spanish Electoral Archive. SEA database}. Scientific Data, 8(1), 160.

\bibitem{piketty2018} 
Piketty, T. (2018). \textit{Brahmin Left vs Merchant Right: Rising Inequality and the Changing Structure of Political Conflict}. WID.world Working Paper.

\bibitem{piketty2021} 
Piketty, T. (2021). \textit{Capital e ideología}. Editorial Planeta / Deusto.

\bibitem{rodriguezpose2018} 
Rodríguez-Pose, A. (2018). \textit{The revenge of the places that don't matter (and what to do about it)}. Cambridge Journal of Regions, Economy and Society, 11(1), 189-209.

\bibitem{sanchez2023} 
Sánchez-García, J., \& Rodon, T. (2023). \textit{Economic perceptions and voting behavior in complex party systems}. Electoral Studies, 82, 102588.

\bibitem{alcaide}
Alcaide Lara, A. (s.f.). \textit{Evolución del voto y narrativas de crisis}.

\bibitem{aguilar2015} 
Aguilar López, R., \& Aquino López, M. A. (2015). \textit{Modelización de la intención de voto con datos composicionales}. REES: Revista de Estadística, Econometría y Sistemas.

\bibitem{lago2013}
Lago, I., \& Martínez, F. (2013). \textit{La penalización del voto disperso en el sistema electoral español}.

\bibitem{lastras2024}
Lastras Rodríguez, A. (2024). \textit{Captura de patrones complejos de bienestar urbano y comportamiento electoral}.

\bibitem{leal2025}
Leal Varón, et al. (2025). \textit{Mapeando el espectro ideológico izquierda-derecha con Machine Learning}.

\bibitem{montero2009}
Montero, J. R., \& Riera, P. (2009). \textit{El umbral natural y el sistema electoral en provincias de magnitud baja}.

\bibitem{stoetzer2019}
Stoetzer, L. F., et al. (2019). \textit{Forecasting elections in multiparty systems: a Bayesian approach}.

\end{thebibliography}
